%% ma: L10-B13-0095
%% lop: 10
%% chuong: 5
%% bai: 13
%% dang: 10.13.12
%% loai: TN4
%% muc: NB
%% dap_an: "A"
%% nguon: "chuyen-de-cac-so-dac-trung-cua-mau-so-lieu-khong-ghep-nhom-toan-10-chuong-trinh-moi.pdf (GV. Phan Nhật Linh) (P03), Bài 02 Dạng 2 Phần I Câu 23"
%% kien_thuc: "Lớp 10 Bài 13 (mốt)"
%% trang_thai: chinh-thuc
%% ghi_chu: "Tự duyệt theo yêu cầu giáo viên 22/10/2026."
\begin{ex}
Bảng sau ghi lại điểm của $40$ học sinh trong một bài kiểm tra môn Toán:
\begin{center}
\begin{tabular}{|c|c|c|c|c|c|c|c|c|}\hline
Điểm & $3$ & $4$ & $5$ & $6$ & $7$ & $8$ & $9$ & $10$\\\hline
Số học sinh & $2$ & $3$ & $7$ & $18$ & $3$ & $2$ & $4$ & $1$\\\hline
\end{tabular}
\end{center}
Mốt của mẫu số liệu là
\choice{\True $M_o=6$}{$M_o=18$}{$M_o=5$}{$M_o=7$}
\loigiai{Giá trị $6$ có tần số lớn nhất ($18$) nên $M_o=6$. Chọn A.}
\end{ex}
